% !Mode:: "TeX:UTF-8"

\begin{cabstract}
随着文化数字化战略的推进，博物馆导览正在从静态讲解的内容提供转向面向学习过程的交互引导。语音讲解器与图文导览虽提升了信息可达性，但在真实参观情境中仍普遍呈现出信息流单向、探索式提问引导难以持续、对参观节奏与空间语境响应不足等问题，使观众的理解往往停留在听到与浏览，难以形成稳定的意义建构与参与推进。近年来，对话式数字人展示被引入博物馆导览场景，提供了自然语言交互与多模态表达的可能，使导览从内容播报走向与观众共同推进。然而，将对话式数字人展示真正嵌入开放场域的导览体验仍面临三类关键挑战：其一，对话容易停留在零散问答与知识检索，缺乏面向导览推进的讲解组织与追问回收策略；其二，相关展示难以持续维护参观过程中动态变化的学习情境与关注点，无法像专业导览员那样依据此刻的场域与观众状态选择讲解时机与内容；其三，多模态表达往往与对话语义、空间对象和引导意图缺乏一致性机制，导致观众体验割裂，进而削弱观众的注意力跟随与现场感。

为应对上述挑战，本文以博物馆学习理论、多模态交互理论与数字人展示研究为基础，提出面向博物馆导览的数字人智能交互设计框架，将面向导览推进的叙事组织机制、面向现场导览的情境编排机制与面向导览感知的外显引导机制整合为一个可运行的闭环机制，从而把讲什么、怎么讲转化为可实现的设计问题。在此基础上，本文在设计方法与系统实现层面开展研究并完成原型验证，主要工作包括：

（1）提出面向博物馆展品知识导览推进的对话组织策略。将展品知识转译为可参考的分层叙事单元，并设计与之配套的追问路径与回收策略，使导览对话能够围绕观众的理解进程逐步展开，在开放式交流中保持节奏与方向，避免停留在碎片化问答与单次检索。

（2）提出面向导览过程的情境设计表征与维护方法。基于博物馆学习理论，将参观中的个人—社会文化—物理情境转化为系统可追踪的情境线索与导览关注点，并形成可解释、可调节的设计表达；同时制定情境在对话推进与空间移动中的更新规则，使系统能够持续把握专业导览的情境依据，从而支撑导览体验的连贯推进。

（3）在叙事组织与情境编排的基础上，提出面向导览感知的外显引导一致性设计机制。构建数字人状态的表达映射，将讲解语义结构与数字人的引导行为联动，使非语言表达与空间对象、讲解意图保持一致，形成可感知的注意力引导，提升导览过程中的引导清晰度与沉浸连贯性。

（4）基于上述框架与机制，本文实现对话式数字人导览原型系统，并通过用户任务实验对原型系统在学习理解、参与感与引导清晰度等维度进行评估。实验结果表明，相较基线方式，基于面向博物馆导览的数字人智能交互设计框架的原型系统能够在保证讲解准确性的同时提升用户的持续跟随与有效追问，并改善整体参观体验评价。本文工作为文化遗产数字传播中的对话式数字人展示提供了可参考的设计框架、可实现的状态机制与多模态一致性路径，并为面向学习情境的对话式数字人展示设计与验证提供参考。
\end{cabstract}

\begin{eabstract}
With the advancement of cultural digitalization, museum guidance is shifting from static content delivery toward interaction-oriented learning support. Although audio guides and graphic-text guides have improved information accessibility, real on-site visits still commonly show one-way information flow, unsustainable follow-up questioning, and insufficient responsiveness to visit pacing and spatial context. As a result, audience understanding often remains at the level of hearing and browsing, rather than developing into stable meaning construction and sustained participation. In recent years, conversational digital human presentation has been introduced into museum guidance, enabling natural language interaction and multimodal expression, and moving guidance from one-way narration to co-progressive interaction with visitors. However, embedding conversational digital human presentation into open-field museum visits still faces three major challenges: dialogue easily degrades into fragmented Q\&A and knowledge retrieval without guidance-oriented organization and recovery strategies; such presentation still struggles to continuously maintain dynamically changing learning context and visitor focus and therefore cannot choose timing and content like professional guides; and multimodal expression often lacks consistency with dialogue semantics, spatial objects, and guidance intent, which fragments user experience and weakens attention following and sense of presence.

To address these challenges, this thesis proposes an intelligent interaction design framework for digital humans in museum guidance, grounded in museum learning theory, multimodal interaction theory, and digital human presentation research. The framework integrates a narrative organization mechanism for guidance progression, a contextual orchestration mechanism for on-site guidance, and an explicit guidance mechanism for guidance perception into an operational closed loop, turning “what to explain” and “how to explain” into implementable design problems. On this basis, the study carries out design-method research, system implementation, and prototype validation. The main contributions are as follows:

(1) A dialogue organization strategy for advancing museum exhibit knowledge guidance. Exhibit knowledge is translated into hierarchical narrative units for reference, with paired follow-up paths and recovery strategies, so that guidance dialogue can progressively unfold along visitors’ understanding process while maintaining pacing and direction in open interaction, instead of staying at fragmented Q\&A or one-shot retrieval.

(2) A contextual design representation and maintenance method for guidance processes. Based on museum learning theory, personal, sociocultural, and physical contexts in visits are transformed into trackable contextual cues and guidance focus points, and represented in an interpretable and adjustable form. Update rules are further defined for dialogue progression and spatial movement, enabling the system to continuously preserve the contextual basis of professional guidance and support coherent guidance progression.

(3) A consistency-oriented explicit guidance design mechanism for conversational digital human presentation in museums. A digital-human state-to-expression mapping is constructed to couple semantic explanation structures with guidance behaviors, ensuring that nonverbal expression remains aligned with spatial objects and explanation intent, thereby forming perceptible attention guidance and improving guidance clarity and immersive continuity.

(4) A conversational digital human guidance prototype system is implemented and evaluated through user task experiments on learning comprehension, engagement, and guidance clarity. Results show that, compared with baseline methods, the prototype system based on the proposed framework improves sustained following behavior and effective follow-up questioning while preserving explanation accuracy, and also improves overall museum guidance experience evaluation. This work provides a design framework for reference, an implementable state mechanism, and a multimodal consistency pathway for conversational digital human presentation in digital cultural heritage communication, and offers practical references for learning-oriented presentation design and validation.
\end{eabstract}
